\section{边疆与财政：每一条道路都有账}
\label{sec:synthesis-frontier-fiscal}

从白登到河西，从西域绿洲到凉州军镇，汉的边疆不是一条固定外线。骑兵、关塞、水源、市场、牧地、使节和地方首领共同组成它；朝廷取得的不是一次永久“占有”，而是必须反复补给、付款与协商的关系。白登后的和亲、关市与边郡守备是在关中恢复、骑兵条件和有限军费下的组合，不应被简单贴作软弱或强硬。

武帝时期的远征把马匹、粮运、边郡动员与新据点的费用推到中枢。卫青、霍去病的军事成功改变朔方和河西的战略关系，也将驻军、移民与水源维持变成日常支出。五铢钱、盐铁、均输与平准使朝廷更容易调度部分资源，不能魔术般消除转运损耗、地方执行和价格风险。《盐铁论》并非逐字会议记录，仍保存了一项难以回避的争执：边费、物价和运输究竟由谁承担。

轮台诏拒绝某些大规模远征与屯田，承认边地劳苦和远距供给有边界，并不等于帝国退回内地。西域都护、屯田和使节继续维系绿洲关系，仍受河西粮道、匈奴竞争与城邦选择限制。赵充国在羌地提出屯田、招抚与分化，也把“进攻还是退守”的抽象二分改成道路、军粮和不同部众的实际选择。

东汉羌地战争、军费、迁徙与招募不断把边地成本压到凉州、并州和地方社会。州郡军事化不只是宫廷失败的结果，也是持续成本无法由稳定的共同财政承担时，地方军队在解决眼前问题中形成自身供给与忠诚网络的结果。帝纪中的远征、简牍中的军粮、钱币的流通痕迹和道路城址必须一同阅读；没有一类材料能给出帝国的完整账目。

\begin{causalchain}[从远征到地方军权]
远征与驻防需要持续的粮、马、道路和中介；财政技术提高了中枢调度的可能，也将执行风险转给地方；当中央无法稳定承担边费，州郡与边将必须自筹兵粮并取得谈判空间。后来的地方军权由此有了物质条件，并非由一纸任命单独造成。
\end{causalchain}

\begin{sourcenote}[边防数字和路线]
《史记》《汉书》的匈奴传、西域传、食货志与《后汉书·西羌传》维持年序和朝廷的解释；居延、肩水金关、悬泉置简牍补足特定关塞的供给细节。兵数、伤亡、精确路线和“平定”范围必须保留文本限度，不能由战报或局部文书外推。
\end{sourcenote}
\subsection{后期的回声}

东汉末年，边费、州郡军政与皇帝名义被不同集团分别使用，提示财政压力并非只在国库中结算。灵帝朝的州郡调整与刘焉的地方经营，是地方权力取得制度入口的不同条件；曹操奉迎献帝又显示，名义须与屯田、军队和道路结合才会成为资源。《后汉书·灵帝纪》《后汉书·刘焉传》《三国志·武帝纪》《后汉书·献帝纪》可说明这些节点，不能将它们串成单一、必然的崩解故事。